\chapter{Conclusions \& Future Work}
\label{ch:conclusions}

\section{Achievements against objectives}
\label{sec:conc_objectives}

All five objectives of Section~\ref{sec:intro_objectives} were met.
\textbf{O1}: the trading MDP was formulated with a complete 9-feature
observable state, a divisible-asset Discrete(3) action model with
feasibility masking, and a fee-inclusive wealth-change reward whose
telescoping property ($\sum r_t = W_T - W_0$) was verified numerically in
every episode. \textbf{O2}: both learning rules were derived from first
principles --- including the likelihood-ratio derivation explaining the
$\log\pi_\theta$ term --- and implemented from scratch in PyTorch with no
RL library in the training path. \textbf{O3}: both implementations were
validated on CartPole, Flappy Bird and LunarLander, clearing random
floors on all three. \textbf{O4}: both were evaluated on 26 held-out
months of SPY data against do-nothing, random and fully invested
buy-and-hold baselines, under a chronological split, matched 80k-step
budgets and three seeds; a PPO row and a three-level fee sensitivity
study completed the ablation. \textbf{O5}: the observed behaviours were
analysed mechanistically, and the limitations below are stated openly.

\section{Principal findings}
\label{sec:conc_findings}

\begin{enumerate}
  \item \textbf{Under a risk-neutral wealth reward, the rational policy
  is buy-and-hold --- and REINFORCE finds exactly that.} In the pilot
  (minimal state) three different policy-gradient optimisers produced
  action sequences identical to fully invested buy-and-hold on all 26
  test months; with the complete 9-feature state, REINFORCE still
  converged to the identical-to-the-cent buy-and-hold policy. The
  collapse is a property of the \emph{objective}, not of any optimiser
  or of the state alone. This is the empirical face of an old
  theoretical point \cite{moody2001}: the reward function defines the
  incentive, and an agent can only be blamed for optimising what it was
  asked to optimise.
  \item \textbf{The algorithm families respond differently to the
  same incentive.} Deep Q-learning, trained on the same MDP, budget and
  data, learned a near-fully-invested policy that trades around local
  turbulence: lower monthly volatility (274 vs.\ 310), lower maximum
  drawdown (4.8\% vs.\ 6.0\%) and higher monthly Sharpe (0.45 vs.\ 0.40)
  than buy-and-hold, at a mean cost of \$1.96/month --- and it adapted
  its trading intensity to transaction fees (19.8 trades/month at zero
  fee, 6.1 at 20 bps). PPO, whose clipping deliberately limits policy
  movement, was bimodal across seeds: exact buy-and-hold on two seeds, a
  cautious low-exposure policy on the third.
  \item \textbf{Evaluation design is a result in itself.} Without the
  properly capitalised buy-and-hold baseline, every learned agent in this
  dissertation would have been reported as beating the market by an order
  of magnitude --- the precise failure mode the financial machine-learning
  methodology literature warns of \cite{bailey2014, lopezdeprado2018}.
  With it, the honest conclusion is recoverable.
\end{enumerate}

\section{Critical appraisal and limitations}
\label{sec:conc_appraisal}

\paragraph{What is limited.} (i) A single asset and a single market
regime window; conclusions about the 2022--2024 test period need not
transfer to other assets or eras. (ii) Daily bars in monthly episodes
rather than the intraday $T{=}390$ horizon originally envisaged; the
environment API is horizon-agnostic, but the empirical claims hold only
at the tested granularity. (iii) Three seeds bound, but do not
eliminate, run-to-run variance --- and the PPO bimodality shows that seed
sensitivity is a real phenomenon here, not a formality. (iv) Episodes
are treated as independent although consecutive months are correlated
(non-IID); the MDP framework used here has no mechanism for that
temporal structure. (v) The state, however deliberately designed,
remains a partial observation of the market: the problem is strictly a
POMDP \cite{kaelbling1998}, and no belief-state or recurrent machinery
was used. (vi) Execution is idealised (fills at the close, unlimited
liquidity at the quoted price); for SPY at these trade sizes the
idealisation is mild, but it is an idealisation.

\paragraph{What could have been done differently.} With hindsight, two
choices would change. First, the pilot's whole-share action model was a
foreseeable dead end --- the divisibility problem it created (fixed
actions becoming infeasible as prices moved) consumed implementation time
that the final divisible-slice model dissolved immediately. Second, the
risk-neutral reward could have been paired from the outset with a
risk-aware variant as a designed experiment; the project discovered the
right question (``what happens to Deep Q's trimming behaviour when the
reward pays for it?'') only after the ablation results arrived, leaving
it as future work rather than a delivered chapter.

\paragraph{What the project got right, in hindsight.} Validating
implementations on games before markets caught errors where they were
diagnosable; the fully invested baseline converted a would-be false
positive into the central finding; fee-inclusive rewards made the fee
sensitivity study possible at all; and deriving the losses before
implementing them was, per the supervision process, the difference
between using an algorithm and understanding one.

\section{Future work}
\label{sec:conc_future}

\begin{itemize}
  \item \textbf{Risk-aware rewards.} Replacing $r_t = \Delta W_t$ with a
        drawdown-penalised variant or the differential Sharpe ratio
        \cite{moody2001} changes the incentive that currently makes
        buy-and-hold optimal. Section~\ref{sec:res_risk} reports one point
        on that frontier and shows the shape of the difficulty: the
        penalty weight controls a trade-off that collapses quickly, with
        abstention rather than better timing as the agent's easiest reply.
        The natural continuation is a reward that penalises variability
        without rewarding inactivity, such as the differential Sharpe ratio
        \cite{moody2001}, evaluated across the frontier rather than at a
        single weight.
  \item \textbf{Intraday horizons.} Minute-bar episodes ($T{=}390$) via
        the same environment API, testing whether the family-level
        behavioural differences persist at higher frequency, where
        fee pressure and state dynamics are sharper.
  \item \textbf{Partial observability.} Recurrent policies, or an
        explicit regime model in the spirit of hidden Markov models
        \cite{rabiner1989}, to address the POMDP structure acknowledged
        throughout.
  \item \textbf{Continuous actions and portfolios.} Replacing the fixed
        trade slice with a continuous allocation, and extending to
        multi-asset weight vectors with risk-adjusted objectives
        \cite{markowitz1952, jiang2017}.
  \item \textbf{Value-based refinements.} Double Q-learning
        \cite{vanhasselt2016} to control the overestimation bias that
        masked TD targets partially mitigate here.
\end{itemize}

\section{Reflection on project management}
\label{sec:conc_pm}

The project followed a staged plan agreed in supervision, summarised in
Table~\ref{tab:conc_plan}. Two properties of the plan proved decisive.
First, \emph{stage gates with demonstrable artefacts}: no stage began
until the previous one produced something checkable (passing tests,
learning curves, result tables), which localised every fault to the stage
that introduced it. Second, \emph{descoping as a deliberate act}: the
most consequential replanning decision was narrowing an early
multi-component ambition (uncertainty-aware portfolio management with a
forecasting subsystem) to a single, fully understood trading MDP with two
from-scratch algorithms. The decision was made with the supervisor,
explicitly trading breadth of feature count for depth of understanding
--- a trade this dissertation's examinable content vindicates.

\begin{table}[h]
\centering
\caption{Project stages and their exit artefacts.}
\label{tab:conc_plan}
\small
\begin{tabular}{llp{6.2cm}}
\toprule
\textbf{Period (2026)} & \textbf{Stage} & \textbf{Exit artefact} \\
\midrule
Feb--Apr & Foundations \& literature & working notes; supervised
derivation exercises; scope agreed \\
May--Jun & Phase 0: games & from-scratch REINFORCE/DQN clearing random
floors on three games; suite metrics \\
Jun--Jul & Pilot trading MDP & minimal-state environment + tests; the
buy-and-hold collapse result and fair-baseline lesson \\
late Jul & Supervision checkpoint & recorded feedback; corrected
objective derivation set as homework; final-model scope fixed \\
Aug & Final model \& ablation & 9-feature environment (8/8 tests),
three-algorithm ablation, fee study, all charts \\
Aug--Sep & Write-up & this document, written methodology-first as
advised; abstract last \\
\bottomrule
\end{tabular}
\end{table}

Risk management was handled by construction rather than by contingency
tables: the riskiest claim in any learning project --- ``the
implementation is correct'' --- was retired early by the game-validation
gate; the riskiest scientific claim --- ``the agent beats the market''
--- was constrained by baselines designed to make it hard to assert
falsely. Weekly supervision meetings served as checkpoints, and the final
experiment suite was engineered to be regenerable in minutes on a laptop
CPU, so that every claim in this dissertation can be demonstrated live at
the viva.

\paragraph{Closing summary (plain language).} Two learning programs were
built from nothing, taught to play games, and then set loose on seven
years of real market prices with imaginary money. One concluded that the
best strategy was the simplest: buy immediately and hold --- which, for
the goal it was given, is mathematically the right answer. The other
learned to step partially out of the market when conditions turned rough,
accepting slightly smaller profits for noticeably smaller losses, and
learned to trade less when trading cost more. Neither beat the market,
and the project was designed so that this honest answer would be
detectable --- because in this field, the experiments that claim
otherwise are usually the ones with the mistake.
